% !TeX spellcheck = en_US
\documentclass[french]{yLectureNote}

\title{Introduction à la termodynamique}
\subtitle{Thermodynamique}
\author{Paulhenry Saux}
\date{\today}
\yLanguage{Français}

\professor{}

\usepackage{graphicx}%----pour mettre des images
\usepackage[utf8]{inputenc}%---encodage
\usepackage{geometry}%---pour modifier les tailles et mettre a4paper
%\usepackage{awesomebox}%---pour les boites d'exercices, de pbq et de croquis ---d\'esactiv\'e pour les TP de PC
\usepackage{tikz}%---pour deiffner + d\'ependance de chemfig
\usepackage{tkz-tab}
\usepackage{tabularx}%---pour dimensionner automatiquement les tableaux avec variable X
\usepackage{awesomebox}%---Pour les boites info, danger et autres
\usepackage{menukeys}%---Pour deiffner les touches de Calculatrice
\usepackage{fancyhdr}%---pour les en-t\^ete personnalis\'ees
\usepackage{blindtext}%---pour les liens
\usepackage{hyperref}%---pour les liens (\`a mettre en dernier)
\usepackage{caption}%---pour la francisation de la l\'egende table vers Tableau
\usepackage{pifont}
\usepackage{array}%---pour les tableaux
\usepackage{lipsum}
\usepackage{yFlatTable}
\usepackage{multicol}
\newcommand{\Lim}[1]{\lim\limits_{\substack{#1}}\:}
\renewcommand{\vec}{\overrightarrow}
\newcommand{\bz}{\overline{z}}
\DeclareMathOperator{\card}{card}
\renewcommand{\vec}{\overrightarrow}
\newcommand{\N}[0]{\mathbb{N}}
\newcommand{\R}[0]{\mathbb{R}}
\newcommand{\C}[0]{\mathbb{C}}
\newcommand{\dd}[0]{\mathrm{d}}
\newcommand{\er}{\vec{e_{\rho}}}
\newcommand{\ep}{\vec{e_{\varphi}}}
\newcommand{\pip}{\pi^+}
\newcommand{\pim}{\pi^-}
\newcommand{\mv}{\mathcal{V}}
\DeclareMathOperator{\Div}{Div}
\newcommand{\rot}{\vec{\mathrm{rot}}}
\newcommand{\grad}{\vec{\mathrm{grad}}}
\begin{document}
\setcounter{chapter}{1}
\chapter{Travail des forces de pression}

% \section{Énergie reçue par un système fermé}

% \explanation{1}{Le travail est une grandeur de transfert : c'est l'intégrale du flux de puissance à travers la frontière du système.}

% \begin{flalign*}
%     W_{i\to f} &= \int_{t_i}^{t_f} \phi(t)\dd t\explain{1}{right}{0}{0.5}{}
% \end{flalign*}

% \explanation{2}{La puissance $\phi$ est l'intégrale de la densité de flux $\varphi$ (en $W.m^{-2}$) sur toute la surface limite $\partial \Omega$.}

% \begin{flalign*}
%     \phi(t) &=  \int_{\partial \Omega} \varphi(M, t)\dd S \explain{2}{right}{0}{0.5}{}
% \end{flalign*}

% \subsection{Point de difficulté : Passage du local au global}
% Le travail élémentaire sur une facette $dS$ pendant $dt$ est :
% \begin{flalign*}
%     \delta^2 W &= \dd\vec{F}(M,t)\cdot \dd \vec{l}(M,t)\\
%     \varphi(M,t) &= \vec{\tau}(M,t)\cdot \vec{v_m}(M,t) \quad \text{(Densité de puissance)}
% \end{flalign*}



\section{Cas du cylindre indéformable avec piston mobile}

\warningInfo{Convention de signe}{L'énergie reçue par le système (le fluide) est comptée positivement ($W > 0$). Si le système fournit du travail, $W < 0$.}

\explanation{a1}{Seule la surface du piston $S_p$ se déplace ($\vec{v_m} \neq \vec{0}$). Sur les parois fixes, le produit scalaire est nul.}

\begin{flalign*}
    W &= \int_{t_i}^{t_f}\int_{S_p} \vec{\tau(t)}\cdot \vec{v_p}(t)\dd S\dd t \explain{a1}{right}{0}{0.5}{}\\
    &= \int_{t_i}^{t_f}\vec{F_{P/\Omega}}\cdot \vec{v_p}(t)\dd t
\end{flalign*}

\subsection{Expression de la force $F_{P/\Omega}$ (via le PFD)}
On étudie le piston (masse $m_p$) pour trouver la force qu'il exerce sur le fluide :
\begin{flalign*}
    \vec{F_{P/\Omega}} = \vec{F_{vol}}  + \vec{F_{atm}} + \vec{F_{op}} - m_p \frac{\dd \vec{v_p}}{\dd t}
\end{flalign*}

\explanation{a2}{L'énergie cinétique du piston est négligée entre l'état initial et final (vitesses nulles), ou la masse du piston est considérée comme constante sans apport d'énergie cinétique externe.}

\explanation{a3}{$P_{ext}$ est une pression fictive "extérieure" qui modélise toutes les forces s'opposant à l'expansion du gaz (poids, atmosphère, opérateur).}

\begin{flalign*}
    W &= \int_i^f -(P_{atm} + \frac{m_p g}{S_p} + \frac{F_{op}}{S_p}) \dd V \explain{a2}{right}{0}{0.5}{}\\
    W &= \int_i^f -P_{ext} \dd V \explain{a3}{left}{0}{0.5}{}
\end{flalign*}
\begin{theorem}[Travail pour un piston]
    $$W =\int_{V_i}^{V_f} -P_{ext} \dd V$$
\end{theorem}
\criticalInfo{Remarque}{
Même si $V_i = V_f$, le travail n'est pas forcément nul si le chemin parcouru implique des variations de volume intermédiaires (le travail dépend du chemin suivi).
}

\section{Applications classiques}

\subsection{Compression monotherme brusque (Irréversible)}
On impose brutalement une pression finale $P_f$. La pression extérieure $P_{ext}$ est constante et égale à $P_f$ pendant toute la transformation.
\begin{flalign*}
    W_{brutal} &= -P_f \int_{V_i}^{V_f} \dd V = -P_f(V_f - V_i)
\end{flalign*}



\subsection{Transformation Quasi-statique (Réversible)}
Le système est à l'équilibre à chaque instant.

\explanation{a5}{À chaque étape infinitésimale, la pression du gaz $P$ est égale à la pression extérieure $P_{ext}$.}
\begin{flalign*}
    W_{qs} &= \int_i^f -P \dd V \explain{a5}{right}{0}{0.5}{}\\
    \text{Pour un GP isotherme : } W_{qs} &= -\int_i^f \frac{nRT}{V}\dd V = nRT \ln\left(\frac{V_i}{V_f}\right)
\end{flalign*}

\subsection{Difficultés}
\begin{itemize}
    \item \textbf{P vs $P_{ext}$ :} On utilise toujours $P_{ext}$ pour le travail. On ne remplace par $P$ (pression du gaz) que si la transformation est quasi-statique.
    \item \textbf{Signe :} Si $V$ diminue ($dV < 0$), alors $-P_{ext}dV > 0$, le travail est bien reçu.
    \item \textbf{Calcul numérique :} Toujours convertir les cm en m et les $P$ en Pascal ($Pa$) avant de calculer $W$ en Joules ($J$).
\end{itemize}
\end{document}


